\documentclass{article}
\usepackage[utf8]{inputenc}
\usepackage{amsmath,amssymb}
\usepackage[most]{tcolorbox}
\usepackage{geometry}
\geometry{margin=2.5cm}

\newcommand{\step}[1]{\par\noindent\hangindent=3.5em\hangafter=1%
  \makebox[3.5em][l]{\textbf{#1.}}}
\newcommand{\steptag}[1]{\hfill{\small\textsf{[#1]}}}

\newtcolorbox{proofbox}[1]{
  colback=gray!5, colframe=gray!60,
  fonttitle=\bfseries, title={#1},
  breakable, enhanced,
  left=4pt, right=4pt, top=4pt, bottom=4pt,
  before skip=10pt, after skip=10pt
}

\begin{document}

\begin{proofbox}{Parametric halo collapse: for an exact signed idempotent ...}
\step{1} Parametric halo collapse: for an exact signed idempotent P with 0 < delta(P) <= 1/4, nonempty visible set W(P), hidden top vertex v of height H, and any halo width a > 0, writing sigma\_a for the positive coefficient mass v places on rows at ell-1 distance > a*tau from conv W (tau = sqrt(delta)), sigma for the invisible mass, and nu\_v for the row negative mass, one has H*(1 - sigma\_a) <= (sigma - sigma\_a)*a*tau + nu\_v*(2 + 4*delta). \steptag{validated} \\[6pt]
\step{1.1} Setup and notation: let alpha\_j=P\_vj, alpha\_j\textasciicircum{}+=max(alpha\_j,0), alpha\_j\textasciicircum{}-=max(-alpha\_j,0), C=conv\{p\_w:w in W(P)\}, d\_j=dist\_1(p\_j,C), tau=sqrt(delta), and A=\{j:d\_j>a*tau\}. Then d\_v=H, d\_j<=H for every row j, sigma\_a=sum\_\{j in A\} alpha\_j\textasciicircum{}+, sigma=sum\_\{d\_j>0\} alpha\_j\textasciicircum{}+, nu\_v=sum\_j alpha\_j\textasciicircum{}-, 0<=sigma\_a<=sigma, and every x in C is within 2+4*delta in l1 of every row p\_k. \steptag{validated} \\[6pt]
\step{1.1.1} From def-height and def-visible-set, C is the visible-row hull, d\_v=dist\_1(p\_v,C)=H for the hidden top vertex v, and d\_j<=H for every row j because H is the maximum row distance from C. \steptag{validated} \\[4pt]
\step{1.1.2} From def-invisible-mass and the definition of A=\{j:d\_j>a*tau\}, sigma=sum\_\{d\_j>0\} alpha\_j\textasciicircum{}+ and sigma\_a=sum\_\{j in A\} alpha\_j\textasciicircum{}+ with A subset \{j:d\_j>0\}; hence 0<=sigma\_a<=sigma, sigma-sigma\_a is the positive mass on rows with 0<d\_j<=a*tau, and nu\_v=sum\_j alpha\_j\textasciicircum{}- by def-negative-mass for row v. \steptag{validated} \\[4pt]
\step{1.1.3} From the row-geometry clause of def-signed-idempotent, any two rows have l1 distance at most 2+4*delta; since C is the convex hull of visible rows, every x in C is a convex combination of rows and therefore ||x-p\_k||\_1<=2+4*delta for every row p\_k. \steptag{validated} \\[4pt]
\step{1.2} Case sigma\_a>=1: the desired inequality holds because H>=0 makes H*(1-sigma\_a)<=0, while sigma-sigma\_a>=0, a*tau>0, nu\_v>=0, and 2+4*delta>=0 make the right-hand side nonnegative. \steptag{validated} \\[4pt]
\step{1.3} Case sigma\_a<1, row-reproduction split: define q=(sum\_\{j notin A\} alpha\_j\textasciicircum{}+ p\_j - sum\_k alpha\_k\textasciicircum{}- p\_k)/(1-sigma\_a). Then 1-sigma\_a>0 and p\_v=sum\_\{j in A\} alpha\_j\textasciicircum{}+ p\_j + (1-sigma\_a)*q. \steptag{validated} \\[6pt]
\step{1.3.1} By def-signed-idempotent, P\textasciicircum{}2=P gives row reproduction p\_v=sum\_j alpha\_j p\_j=sum\_j alpha\_j\textasciicircum{}+ p\_j - sum\_k alpha\_k\textasciicircum{}- p\_k. By lem-mass-split, sum\_j alpha\_j\textasciicircum{}+=1+nu\_v=sum\_j alpha\_k\textasciicircum{}-+1, so sum\_\{j notin A\} alpha\_j\textasciicircum{}+ - sum\_k alpha\_k\textasciicircum{}- = 1-sigma\_a; since this case assumes sigma\_a<1, substituting the displayed q gives the asserted split. \steptag{validated} \\[4pt]
\step{1.4} Case sigma\_a<1, lower residual bound: applying lem-residual-lower to C, p=p\_v, coefficients alpha\_j\textasciicircum{}+ on j in A, and the q from the row-reproduction split gives H<=dist\_1(q,C). \steptag{validated} \\[6pt]
\step{1.4.1} The split from node 1.3 has p\_v=sum\_\{j in A\} alpha\_j\textasciicircum{}+ p\_j +(1-sigma\_a)q with coefficients alpha\_j\textasciicircum{}+>=0 and s=sum\_\{j in A\} alpha\_j\textasciicircum{}+=sigma\_a<1. For each j in A, node 1.1 gives dist\_1(p\_j,C)=d\_j<=H=dist\_1(p\_v,C). Hence lem-residual-lower applies (with the empty-A case giving q=p\_v directly) and yields H<=dist\_1(q,C). \steptag{validated} \\[4pt]
\step{1.5} Case sigma\_a<1, upper residual bound: applying lem-residual-upper to the same q with positive weights alpha\_j\textasciicircum{}+ for j notin A, negative weights alpha\_k\textasciicircum{}- for all negative coefficients, and D\_k=2+4*delta gives (1-sigma\_a)*dist\_1(q,C)<=sum\_\{j notin A\} alpha\_j\textasciicircum{}+ d\_j + nu\_v*(2+4*delta). \steptag{validated} \\[6pt]
\step{1.5.1} For lem-residual-upper set b\_j=alpha\_j\textasciicircum{}+ and p\_j=p\_j for indices j notin A, set c\_k=alpha\_k\textasciicircum{}- and r\_k=p\_k for all negative coefficients, and set m=sum\_\{j notin A\}alpha\_j\textasciicircum{}+-sum\_k alpha\_k\textasciicircum{}-=1-sigma\_a>0 by node 1.3. The displayed formula for q is exactly (sum b\_j p\_j - sum c\_k r\_k)/m. Node 1.1 gives ||x-r\_k||\_1<=2+4*delta for all x in C, so D\_k=2+4*delta is admissible; since sum\_k c\_k=nu\_v, lem-residual-upper yields the asserted bound. \steptag{validated} \\[4pt]
\step{1.6} Halo pricing: sum\_\{j notin A\} alpha\_j\textasciicircum{}+ d\_j <= (sigma-sigma\_a)*a*tau, because the d\_j=0 terms contribute zero and the remaining j notin A with d\_j>0 have d\_j<=a*tau and total positive mass sigma-sigma\_a. \steptag{validated} \\[4pt]
\step{1.7} Conclusion: in the case sigma\_a<1, multiply H<=dist\_1(q,C) by the positive factor 1-sigma\_a and combine the upper residual bound and halo pricing to obtain H*(1-sigma\_a) <= (sigma-sigma\_a)*a*tau + nu\_v*(2+4*delta); with the sigma\_a>=1 case this proves the root claim. \steptag{validated} \\[4pt]
\end{proofbox}

\end{document}
